% Suggested time: 1-2 minutes
\section{Conclusion}

\begin{frame}{Interpreting the Route}
  \small
  \begin{itemize}
    \item Meets the objective's constraints: leave USTP, visit each of the 9 landmarks \textbf{exactly once},
          return to USTP, in one loop of \textbf{\DCTotal{} km}.
    \item Nearby places are grouped: the Divisoria cluster V4, V5, V6 is visited consecutively
          (0.54 and 0.74 km).
    \item Driven with the raw one-way distances, the same route is \textbf{13.54 km}; 13.82 km uses the
          averaged, undirected weights of the model.
  \end{itemize}
  \medskip
  \begin{block}{Where the route loses to the optimum}
    \footnotesize
    Only the order inside the cluster differs: ours V3$\to$V4$\to$V5$\to$V6 ($1.31+0.54+0.74 = 2.59$ km),
    optimal V3$\to$V5$\to$V4$\to$V6 ($0.97+0.54+0.87 = 2.38$ km). Merge M2 must add the shortest bridge
    V6--V5, so the edge V6--V4 is never considered. The long leg V8$\to$V0 is \emph{not} the cause:
    the optimal tour uses it too.
  \end{block}
\end{frame}

\begin{frame}{Is It Optimal? No --- and We Can Show It}
  \small
  Exact Held--Karp~\cite{heldkarp1962} on the same graph, used only for comparison:
  \begin{center}
    \OptRoute{} \quad = \textbf{\OptTotal{} km}
  \end{center}
  Our Divide and Conquer route: \textbf{\DCTotal{} km}, i.e.\ \textbf{\DCGap\% above optimal}, so it is an
  \textbf{approximate} solution.
  \medskip

  \centering
  \scriptsize
  % AUTO-GENERATED by code/main.py --export-latex. Do not edit; re-run the export instead.
\begin{tabular}{l r r r r r r r}
\toprule
\textbf{Method} & \textbf{Best} & \textbf{Worst} & \textbf{Mean} & \textbf{Std.} & \textbf{Mean gap} & \textbf{Optimal} & \textbf{ms/run} \\
\midrule
Shortest bridge (paper) & 13.61 & 14.26 & 13.88 & 0.25 & 2.01\% & 6/100 & 0.405 \\
Best splice (variant) & 13.61 & 14.13 & 13.91 & 0.24 & 2.20\% & 17/100 & 0.342 \\
Held--Karp (exact) & 13.61 & 13.61 & 13.61 & 0.00 & 0.00\% & -- & -- \\
\bottomrule
\end{tabular}


  \medskip
  \footnotesize The paper reports 10 runs; we repeated ours over \RunCount{} random seeds. Only \BridgeOptHits{}
  runs happened to be optimal, and the algorithm cannot tell which; only Held--Karp shows it.
\end{frame}

\begin{frame}{Limitations and Improvements}
  \begin{columns}[T]
    \begin{column}{0.48\textwidth}
      \begin{block}{Limitations}
        \small
        \begin{itemize}
          \item Approximate (\DCGap\% above optimal); depends on the seed (\BridgeBest{}--\BridgeWorst{} km
                over \RunCount{} seeds)
          \item Distance only: no traffic or time windows
          \item One-way streets averaged (13.54 km on raw one-way values)
          \item OSRM vs.\ Google Maps: 13.82 vs.\ 14.40 km; optimality is shown for the OSRM matrix only
        \end{itemize}
      \end{block}
    \end{column}
    \begin{column}{0.48\textwidth}
      \begin{block}{Improvements}
        \small
        \begin{itemize}
          \item 2-opt: reversing V4$\to$V5 swaps V3--V4, V5--V6 ($1.31+0.74$) for V3--V5, V4--V6
                ($0.97+0.87$), saving 0.21 km $\Rightarrow$ exactly \OptTotal{} km
          \item Restart with several seeds, keep the best (best of \RunCount{}: \BridgeBest{} km; only
                Held--Karp shows it is optimal)
          \item Best-splice merge: more optimal hits, but no better on average
          \item Travel-time weights from Google Maps
        \end{itemize}
      \end{block}
    \end{column}
  \end{columns}
\end{frame}

\begin{frame}{Conclusion}
  \begin{itemize}
    \item Modelled a downtown Cagayan de Oro tour as a TSP on $K_{\NumVertices}$ with road distances (km).
    \item Implemented Divide and Conquer (Algorithm 2.8), fixing the merge so that it always gives a
          valid Hamiltonian cycle, in $O(n^2)$ time.
    \item Route: \DCRoute{}
    \item Total \textbf{\DCTotal{} km}, verified; \textbf{\DCGap\% above} the optimum of \OptTotal{} km.
    \item Fast and valid, but \textbf{approximate}, not proven optimal.
  \end{itemize}
\end{frame}

\begin{frame}[allowframebreaks]{References}
  \tiny
  \nocite{lunagarcia2024performance, karp1977probabilistic, codingdrills_dc_tsp, morenomaza_dc, osm, googlemaps, luxen2011osrm, hunter2007matplotlib, applegate2006tsp}
  \bibliographystyle{plain}
  \bibliography{dc-tsp-refs}
\end{frame}
